% !TeX program = lualatex
% Definition and sources only; this file is not an LLM solution attempt.
\documentclass[11pt,a4paper]{article}
\usepackage[margin=25mm]{geometry}
\usepackage{fontspec}
\setmainfont{lmroman10-regular.otf}[BoldFont=lmroman10-bold.otf,
 ItalicFont=lmroman10-italic.otf,BoldItalicFont=lmroman10-bolditalic.otf]
\usepackage{amsmath,amssymb,mathtools}
\usepackage{unicode-math}
\setmathfont{latinmodern-math.otf}
\AtBeginDocument{\let\setminus\smallsetminus}
\usepackage{xurl,hyperref}
\hypersetup{colorlinks=true,urlcolor=blue}
\setcounter{secnumdepth}{0}
\setlength{\parindent}{0pt}
\setlength{\parskip}{5pt}
\setlength{\emergencystretch}{3em}
\begin{document}
\section*{If F is a finite field with at least 4 elements and A is an invertible n by n matrix over F, then there are vectors x, y}\label{statement}

\subsection{Definitions and mathematical statement}
Let $\mathbb F_q$ be a finite field of prime-power order $q\ge4$,
let $n\ge1$, and let $A\in\operatorname{GL}_n(\mathbb F_q)$ be
an invertible matrix. A vector is nowhere zero if each of its
coordinates belongs to $\mathbb F_q^*=\mathbb F_q\setminus\{0\}$.
Jaeger's question asks whether there always exists
\[
              x\in(\mathbb F_q^*)^n
                     \quad\text{with}\quad
              Ax\in(\mathbb F_q^*)^n.
\]
Equivalently, must the two sets
$(\mathbb F_q^*)^n$ and $A((\mathbb F_q^*)^n)$ intersect?
In row form, the problem is to choose all $x_j\ne0$ so that
$\sum_{j=1}^n A_{ij}x_j\ne0$ simultaneously for all rows $i$.

All arithmetic takes place in the field, so the statement includes
nonprime orders such as $q=4$. The vectors are allowed to depend
on the matrix, and no lower bound on the number of eligible
vectors is required. The hypotheses of invertibility and the
lower bound on $q$ are part of the assertion; merely having no
zero row is not the same assumption. A vector with some nonzero
coordinate is insufficient: both it and its image must avoid
zero in every coordinate. No requirement is placed on the
individual entries of the invertible matrix itself.
\subsection{Short English statement}
Over every finite field with at least four elements, does every invertible matrix map some nowhere-zero vector to another nowhere-zero vector?
\subsection{Sources}
\begin{itemize}
\item[S1] ULAM.AI and the UnsolvedMath Contributors, supplied problems.json, catalogue entry, AMR Open Problem Lists. Catalogue identity and source transcription; CC BY 4.0.
\url{https://huggingface.co/datasets/ulamai/UnsolvedMath}.
\item[S2] Cooper - Combinatorial Problems I Like (2020 snapshot), Combinatorial Number Theory, source-order bullet 60. 
\url{https://people.math.sc.edu/cooper/combprob.html}.
\end{itemize}
\end{document}
